\section{Introduction}

Brain tumors exhibit substantial variability in morphology, anatomical location, and malignancy, making accurate MRI-based diagnosis challenging. Conventional radiological assessment can suffer from high misdiagnosis rates because tumor appearance varies across patients and tumor types~\cite{yan2016mri}. Recent deep learning methods have shown strong performance for brain tumor classification and segmentation, but many models remain difficult to interpret. This limits their usefulness in clinical decision support, where clinicians need not only a prediction but also evidence showing where the model is looking and whether the decision is consistent with visible tumor regions.

We propose an explainable multi-stage deep learning framework for brain tumor MRI analysis. The pipeline first uses ResUNet to segment the tumor region, then uses the predicted mask to extract an ROI crop. To reduce information loss from cropping, the final classifier uses a dual-input ResNet50 design that combines the ROI crop with the original full MRI image. The system further provides visual explanations using Grad-CAM and generates structured clinical summaries using a MedGemma-based VLM with a deterministic fallback template. This design aims to balance predictive performance with interpretability by linking the final classification decision to explicit localization, visual attention, and textual explanation.

Using the Kaggle Brain Tumor MRI Dataset and the SciDB mask-annotated dataset, our experiments show that ResUNet achieves the best segmentation overlap among the evaluated segmentation backbones, reaching a Dice score of 0.7431 and IoU of 0.6558. For classification, the dual-input ResNet50 achieves 99.00\% accuracy and 98.95\% F1-score, approaching the full-image baseline while preserving the localization benefits of segmentation-guided classification. Ablation experiments show that increasing the ROI margin improves cropped-only classification, but the dual-input design is more effective because it restores global anatomical context. VLM comparisons further show that MedGemma provides the best task-specific summary quality among the evaluated summary-generation models.